% !TEX root = ../main.tex
\chapter{付録B　「手を動かす」の解答例と期待出力}

\section{解答例の読み方}

\term{解答例}は、唯一の書き方ではなく、入力、処理、出力を確認する基準です。
実行結果が異なる場合は、OS、言語版、入力、時刻、通信先を記録します。
各章の演習は、答えを写すことより、観測結果を説明できることを到達条件にします。

\section{第1章から第4章}

第1章では、10進数、2進数、16進数の値が一致することを確認します。
同じビット列を符号なし整数、符号付き整数、文字として読むと、意味が変わることを説明します。

第2章では、式をトークン列とASTへ変換し、後順走査で評価します。
第3章では、配列、スライス、マップの操作時間とメモリ共有を観測します。
第4章では、一つの変更要求が越える境界を数え、責務分割の根拠を示します。

\section{第5章から第8章}

第5章の\code{dig +trace}では、ルート、TLD、権威DNSへ問い合わせ先が移ることを確認します。
第6章の\code{curl -w}では、DNS、接続、TLS、最初の1バイト、全体の時間を分けます。
第7章では、学習データと検証データの損失を比較し、過学習の有無を説明します。
第8章では、Network、Elements、Performanceを同じ操作時刻へ対応づけます。

\section{第9章}

総合演習では、正常系と異常系の観測表を作ります。
症状ごとに、観測点、仮説、反証条件、次の操作を一行ずつ対応づけます。
未確認の原因は断定せず、観測によって候補から外せる形にします。

\section{資料と索引の用語}

\term{索引}は、用語の初出や関連箇所を探すための一覧です。
\term{参考文献}は、本文の根拠や追加学習へ使った資料の一覧です。
\term{DOI}はDigital Object Identifierの略で、論文などを持続的に識別します。
\term{RFC}は、インターネット技術の仕様や知見を公開する文書系列です。
\term{標準仕様}は、実装間で共有する要件や意味を定める文書です。
\term{ライセンス}は、著作物やソフトウェアを利用できる条件です。
\term{リンク切れ}は、記載したURLから資料へ到達できない状態です。
